% !TEX program = lualatex
\documentclass[aspectratio=169,professionalfonts]{beamer}
\newcommand{\slidemode}{1}  % カラー投影モード
\usepackage{beamer_template}
\usepackage{enumerate}

\newcommand{\LectureNum}{4}
\newcommand{\LectureTitle}{逆ラプラス変換}
\newcommand{\CourseName}{応用数学}
\renewcommand{\TermName}{前期}

\begin{document}

\begin{frame}{第4回　逆ラプラス変換}
  \begin{block}{本回の内容}
  逆ラプラス変換
  \end{block}
  \vfill\hfill{\scriptsize \textcolor{gray}{第2章 ラプラス変換　§1}}
\end{frame}

\begin{frame}{定義：逆ラプラス変換}
  \begin{block}{}
  \end{block}
  \vfill\hfill{\scriptsize \textcolor{gray}{p.\,56--61}}
\end{frame}

\begin{frame}{部分分数分解による逆ラプラス変換}
  \begin{block}{}
  \end{block}
\end{frame}

\begin{frame}{例題12}
  \begin{exampleblock}{問題}
    次の逆ラプラス変換を求めよ（ $\alpha, \beta$ は定数）\\[2pt]
    $(1) 1/(s-\alpha)$\\[2pt]
    $(2) 1/(s-\alpha)^{2}$\\[2pt]
    $(3) 1/((s-\alpha)(s-\beta))  \alpha \neq \beta$\\[2pt]
    $(4) 1/((s-\alpha)^{2}+\beta^{2})  \beta \neq 0$\\[2pt]
  \end{exampleblock}
\end{frame}

\begin{frame}{例題12　解答}
  \begin{exampleblock}{解答}
  \end{exampleblock}
\end{frame}

\begin{frame}{例題13}
  \begin{exampleblock}{問題}
    次の関数の逆ラプラス変換を求めよ\\[2pt]
    $(1) 1/((s-1)^{2}+4)$\\[2pt]
    $(2) (s+2)/((s+2)^{2}+9)$\\[2pt]
  \end{exampleblock}
\end{frame}

\begin{frame}{例題13　解答}
  \begin{exampleblock}{解答}
  \begin{enumerate}[1.]
    \item[] 分母を $(s-\alpha)^{2}+\beta^{2}$ の形に整理 $: \alpha=1, \beta=2$
    \item[] $L^{-1}[1/((s-1)^{2}+4)] = L^{-1}[(1/2)\cdot2/((s-1)^{2}+2^{2})]$
    \item[] 基本公式 $L^{-1}[\beta/((s-\alpha)^{2}+\beta^{2})] = e^{\alpha t}\sin(\beta t)$ を適用
    \item[] 分母を $(s-\alpha)^{2}+\beta^{2}$ の形に整理 $: \alpha=-2, \beta=3$
    \item[] 分子 $s+2 = (s-\alpha)|_{\alpha=-2}$ に一致
    \item[] 基本公式 $L^{-1}[(s-\alpha)/((s-\alpha)^{2}+\beta^{2})] = e^{\alpha t}\cos(\beta t)$ を適用
  \end{enumerate}
  \end{exampleblock}
\end{frame}

\begin{frame}{例題14}
  \begin{exampleblock}{問題}
    次の逆ラプラス変換を求めよ\\[2pt]
    (1) 1/((s+2)(s+3)(s+4))\\[2pt]
    $(2) (3s^{2}+5) / (s(s+1)^{2}(s-3))$\\[2pt]
  \end{exampleblock}
\end{frame}

\begin{frame}{例題14　解答}
  \begin{exampleblock}{解答}
  \end{exampleblock}
\end{frame}

\begin{frame}{演習問題（p.\,56--61）}
  \begin{exampleblock}{自分で解いてみよう}
  \begin{description}
    \item[問17] 次の逆ラプラス変換を求めよ $: (1) 1/(s^{2}-2s+5) (2) 1/(s^{2}-4s+8)$
    \item[問13\_演習] 次の関数の逆ラプラス変換を求めよ
    $(1) 1/(s-4)^{2}$\\
    $(2) 1/((s-1)^{2}+4)$\\
    \item[問18] 次の逆ラプラス変換を求めよ
    $(1) (s-1)/(s^{2}-4s+5)$\\
    $(2) (s-2)/(s^{2}-6s+10)$\\
    \item[問19] 次の逆ラプラス変換を求めよ
    (1) 1/(s(s-1)(s-2))\\
    (2) (5s+1)/(s(s-1)(s+2))\\
  \end{description}
  \end{exampleblock}
\end{frame}

\end{document}
